\section{Metoodika} \label{metoodika}

Käesolev peatükk kirjeldab uurimuse metoodikat, sealhulgas AI mudelite valiku kriteeriume, eksperimentide ülesehitust, sisendi ettevalmistamist ning hindamiskriteeriumeid.

\subsection{Mudelite valiku kriteeriumid} \label{subsec:valiku-kriteeriumid}

Uurimuse jaoks mudelite valimisel lähtuti järgmistest kriteeriumidest:

\begin{enumerate}
    \item \textbf{Eesti keele tugi} -- mudel peab suutma genereerida grammatiliselt korrektset eesti keelt ning mõistma eestikeelset sisendit.
    \item \textbf{Kontekstiakna suurus} -- bakalaureusetöö genereerimine eeldab mahuka konteksti (juhendid, näidistööd, varasemalt genereeritud osad) töötlemist, mistõttu on vajalik piisavalt suur kontekstiaken.
    \item \textbf{Kättesaadavus ja dokumentatsioon} -- mudel peab olema kättesaadav API kaudu või veebirakendusena ning omama piisavat dokumentatsiooni.
    \item \textbf{Hind} -- mudeli kasutamine peab olema majanduslikult mõistlik üliõpilase jaoks.
    \item \textbf{Ajakohasus} -- mudel peab olema üks uuemaid versioone (2025--2026).
\end{enumerate}

Nende kriteeriumide alusel valiti neli mudelit: Google Gemini 3.1 Pro, OpenAI GPT-5.4, Anthropic Claude Sonnet 4.6 ning Anthropic Claude Opus 4.6. Opus 4.6 lisati võrdlusse kui sama mudeliperekonna võimekaim mudel, et hinnata mudeli suuruse mõju genereeritud teksti kvaliteedile. Valiku põhjendused on esitatud peatükis~\ref{mudelid}.

\subsection{Eksperimentide ülesehitus} \label{subsec:eksperimentide-ulesehitus}

Eksperimendid viidi läbi kolmes etapis:

\textbf{Etapp 1: Konteksti ettevalmistamine.} Esmalt koostati igale mudelile antav kontekstipõhi, mis sisaldas:
\begin{itemize}
    \item TÜ ATI lõputööde juhendi põhilised nõuded (struktuur, vormistus, hindamiskriteeriumid);
    \item LOTE lõputööde juhendi olulised lõigud;
    \item kolme varasema, hinnetega ``A'' või ``B'' hinnatud bakalaureusetöö sissejuhatused ja kokkuvõtted;
    \item eestikeelse akadeemilise stiili näited ja terminoloogia.
\end{itemize}

\textbf{Etapp 2: Genereerimine.} Iga mudeli puhul genereeriti järgmised lõputöö komponendid:
\begin{itemize}
    \item sissejuhatus (ca 2 lehekülge);
    \item kirjanduse ülevaate peatükk (ca 5--8 lehekülge);
    \item metoodika kirjeldus (ca 3--5 lehekülge);
    \item tulemuste peatükk koos tabelite ja joonistega (ca 5--8 lehekülge);
    \item kokkuvõte (ca 1--2 lehekülge).
\end{itemize}

Genereerimine toimus iteratiivselt: esmalt anti mudelile ülesande kirjeldus koos kontekstiga, seejärel paluti genereerida vastav osa. Vajadusel kasutati järelküsimusi (\textit{follow-up prompts}) teksti täiendamiseks ja parandamiseks.

Lisaks üksikute peatükkide genereerimisele viidi läbi ka tervikliku lõputöö genereerimise eksperiment (eksperiment~6, vt peatükk~\ref{subsec:exp-tervik}), kus igale mudelile anti ülesandeks genereerida täielik eestikeelne bakalaureusetöö. Gemini 3.1 Pro, GPT-5.4 ja Claude Sonnet 4.6 puhul kasutati tavapärast vestlusliidest, samas kui Claude Opus 4.6 puhul kasutati Claude Code'i agentset tööriista (vt peatükk~\ref{claudecode}), mis võimaldas mudelil otse \LaTeX{} faile kirjutada, kompileerida ja hallata. See eksperiment oli kavandatud hindamaks mudelite võimekust järgida keerulisi, mitmeosalisi juhiseid ning säilitada keelelist ja sisulist järjepidevust pikema teksti ulatuses.

\textbf{Etapp 3: Hindamine.} Genereeritud tekstid hindas käesoleva bakalaureusetöö autor mitme kriteeriumi alusel (vt alamlõik~\ref{subsec:hindamiskriteeriumid}).

\subsection{Sisendi ettevalmistamine ja promptide disain} \label{subsec:sisendi-ettevalmistamine}

Promptide (viipade) disain on AI-põhise tekstigenereerimise üks kriitilisemaid aspekte. Serikov ja Biloshchytskyi \cite{serikov2025prompt} on näidanud, et promptimismeetodid mõjutavad oluliselt genereeritud sisu kvaliteeti: \textit{few-shot} õppimine töötas kõige paremini eesmärkide genereerimiseks, samas kui ahelpõhjendamine (\textit{chain-of-thought}) oli efektiivsem professionaalsete standardite jaoks. Kresevic et al. \cite{kresevic2024rag} leidsid, et konteksti struktureeritud vormindamine ja kvaliteetne sisend parandasid GPT-4 Turbo täpsust oluliselt rohkem kui näidete hulga suurendamine, rõhutades andmekvaliteedi ülekaalu andmekoguse ees. Uurimuses kasutati kahte peamist lähenemist:

\textbf{Nullkonteksti viibimine} (\textit{zero-shot prompting}) -- mudelile antakse ülesanne ilma näideteta:
\begin{quote}
\textit{``Kirjuta Tartu Ülikooli arvutiteaduse bakalaureusetöö sissejuhatus teemal X. Sissejuhatus peab olema 2 lehekülge pikk, akadeemilises stiilis ning sisaldama viiteid.''}
\end{quote}

\textbf{Kontekstipõhine viibimine} (\textit{few-shot prompting with context}) -- mudelile antakse koos ülesandega ka juhendid ja näited:
\begin{quote}
\textit{``Järgnevalt on esitatud TÜ ATI lõputööde juhendi nõuded: [juhend]. Siin on näide heast bakalaureusetöö sissejuhatusest: [näide]. Kirjuta sarnases stiilis sissejuhatus teemal X.''}
\end{quote}

Mõlema lähenemise tulemusi võrreldi, et hinnata konteksti mõju genereeritud teksti kvaliteedile.

Lisaks katsetati ka süsteemiviipadega (\textit{system prompts}), mille kaudu anti mudelile roll:
\begin{quote}
\textit{``Sa oled kogenud akadeemiline kirjutaja, kes on spetsialiseerunud eestikeelsete arvutiteaduse lõputööde kirjutamisele. Järgi Tartu Ülikooli vormistusnõudeid.''}
\end{quote}

\subsection{Hindamiskriteeriumid} \label{subsec:hindamiskriteeriumid}

Genereeritud tekstide hindamiseks kasutati järgmist kriteeriumite raamistikku:

\begin{table}[ht]
\centering
\caption{Hindamiskriteeriumide raamistik}
\label{tab:hindamiskriteeriumid}
\begin{tblr}{
    colspec={l X[1.5] X[2]},
    row{1}={font=\bfseries},
    hlines,
    vlines,
    row{even}={bg=colorCellHighlight},
}
Kriteerium & Kirjeldus & Hindamisskaala \\
Struktuur & Teksti loogiline ülesehitus, peatükkide jaotus, sidusus & 1--5 (1=puudub struktuur, 5=suurepärane) \\
Keeleline korrektsus & Grammatika, ortograafia, lauseehitus eesti keeles & 1--5 (1=rohked vead, 5=peaaegu veatu) \\
Tehniline täpsus & Faktide, terminoloogia ja tehnilise sisu õigsus & 1--5 (1=rohked vead, 5=täpne) \\
Akadeemiline stiil & Akadeemilisele tekstile vastav keel ja toon & 1--5 (1=mitteakadeemiline, 5=suurepärane) \\
Viited & Viidete olemasolu, asjakohasus ja korrektsus & 1--5 (1=puuduvad/valed, 5=korrektsed) \\
Praktiline väärtus & Teksti kasutatavus lõputöö alusena & 1--5 (1=kasutuskõlbmatu, 5=valmis kasutamiseks) \\
\end{tblr}
\end{table}

Hindamise viis läbi käesoleva bakalaureusetöö autor (informaatika üliõpilane). Hinnangu stabiilsuse tagamiseks ja ühekordsest lugemisest tuleneva subjektiivsuse vähendamiseks vaatas autor iga genereeritud teksti läbi mitmel korral eri aegadel ning lõplik hinne kujunes nende iteratsioonide põhjal. Ühe hindaja kasutamise piirangut on täpsemalt käsitletud peatükis~\ref{arutelu}.
